\documentclass[11pt,a4paper]{article}
\usepackage[margin=20mm,headheight=15pt]{geometry}
\usepackage{fontspec}
\setmainfont{TeX Gyre Pagella}
\setmonofont{DejaVu Sans Mono}[Scale=0.82]
\usepackage{amsmath}
\usepackage{unicode-math}
\setmathfont{TeX Gyre Pagella Math}
\setcounter{MaxMatrixCols}{20}
\usepackage{microtype}
\usepackage{enumitem}
\usepackage{fvextra}
\usepackage{booktabs}
\usepackage{longtable}
\usepackage{fancyhdr}
\usepackage[hidelinks]{hyperref}
\setlength{\parindent}{0pt}
\setlength{\parskip}{6pt}
\setlength{\emergencystretch}{3em}
\hyphenpenalty=10000
\exhyphenpenalty=10000
\raggedright
\DefineVerbatimEnvironment{Code}{Verbatim}{fontsize=\small,
  breaklines=true,breakanywhere=true,breaksymbolleft={},breaksymbolright={},
  frame=single,framesep=5pt}
\setlist[enumerate]{leftmargin=7mm,itemsep=3pt,topsep=5pt,parsep=0pt}
\pagestyle{fancy}
\fancyhf{}
\fancyhead[L]{\small Computer Vision 1}
\fancyhead[R]{\small Exam 1 · Questions}
\fancyfoot[C]{\small \thepage}
\hypersetup{pdftitle={Exam 1 — Foundations under pressure},pdfsubject={CV1 practice material}}
\begin{document}
{\LARGE\bfseries Exam 1 — Foundations under pressure}\par
{\large Questions only}\par\bigskip
\textbf{Name:} \rule{70mm}{0.4pt}\hfill \textbf{Date:} \rule{30mm}{0.4pt}
\section*{Conditions}
\begin{itemize}[leftmargin=5mm,itemsep=5pt]
\item 50 questions: 30 single-answer multiple choice and 20 numerical. Each is worth 1 raw practice point.
\item Suggested training limit: \textbf{150 minutes}. This is not a claim about the actual exam duration or question count.
\item Use up to three double-sided A4 sheets,\allowbreak{} prepared by you,\allowbreak{} and a local Jupyter notebook. No internet,\allowbreak{} AI assistance or answer key during the attempt.
\item Use Python/\allowbreak{}NumPy for computation; SciPy is also allowed. You do not need executable OpenCV or PyTorch. Specialized routines are provided when needed.
\item For multiple choice,\allowbreak{} select exactly one of A–D. Attempt every question.
\item Unless an exact integer is requested,\allowbreak{} round the final numerical answer to \textbf{two decimal places}. Keep full precision in intermediate steps. Use fractions rather than percentages unless specified.
\item Numerical practice grading uses the syllabus rule: within ±0.01 of the unrounded reference,\allowbreak{} inclusive. Counts require the exact integer; no partial credit.
\item Conventions stated in a question override defaults. Arrays use [row,\allowbreak{}column] =\allowbreak{} [y,\allowbreak{}x]; geometric vectors are columns; image-plane projection is virtual/\allowbreak{}front-facing; foreground is 1.
\item This is original,\allowbreak{} challenging training material based on the supplied course files,\allowbreak{} not a lecturer-issued mock or a prediction of the exam.
\end{itemize}
Solutions are in a separate PDF. Do not open them until you finish.
\clearpage\section*{Part A — Single-answer multiple choice}
\begin{minipage}{\linewidth}
\textbf{Q01 — Multiple choice}\par\smallskip
A camera uses the world-to-camera convention \(X_c = R X_w + t\). Which expression gives its optical centre C in world coordinates?
\begin{enumerate}[label=\textbf{\Alph*.}]
\item \(C=-R t\)
\item \(C=-R^{\mathsf T}t\)
\item \(C=R^{\mathsf T}t\)
\item \(C=t\)
\end{enumerate}
\end{minipage}
\par\vspace{12pt}
\begin{minipage}{\linewidth}
\textbf{Q02 — Multiple choice}\par\smallskip
An image is resized to half its original width and height,\allowbreak{} with no crop. Use the ideal continuous-coordinate model,\allowbreak{} with no pixel-centre offset. Which change is correct?
\begin{enumerate}[label=\textbf{\Alph*.}]
\item Halve \(f_x\),\allowbreak{} \(f_y\),\allowbreak{} skew,\allowbreak{} \(c_x\) and \(c_y\); keep R and t unchanged.
\item Halve only \(f_x\) and \(f_y\); leave the principal point unchanged.
\item Halve R and t; leave K unchanged.
\item Halve every entry of K,\allowbreak{} including \(K_{2,2}\),\allowbreak{} without renormalizing it.
\end{enumerate}
\end{minipage}
\par\vspace{12pt}
\begin{minipage}{\linewidth}
\textbf{Q03 — Multiple choice}\par\smallskip
Under an ideal pinhole model,\allowbreak{} two distinct visible 3D points have exactly the same image coordinates. What must hold?
\begin{enumerate}[label=\textbf{\Alph*.}]
\item They lie on the same viewing ray from the optical centre.
\item They have the same physical size.
\item They have the same world X and Y coordinates,\allowbreak{} regardless of camera pose.
\item They have the same depth.
\end{enumerate}
\end{minipage}
\par\vspace{12pt}
\begin{minipage}{\linewidth}
\textbf{Q04 — Multiple choice}\par\smallskip
Keep illumination,\allowbreak{} material and surface orientation fixed. For an ideal Lambertian patch,\allowbreak{} what is independent of viewing direction?
\begin{enumerate}[label=\textbf{\Alph*.}]
\item Its outgoing radiance,\allowbreak{} apart from visibility and camera effects.
\item Its total projected area in the image.
\item Its brightness under any change of light direction.
\item Its surface normal in every coordinate frame.
\end{enumerate}
\end{minipage}
\par\vspace{12pt}
\begin{minipage}{\linewidth}
\textbf{Q05 — Multiple choice}\par\smallskip
Twenty calibrated light vectors all lie in one plane through the origin. In the linear photometric-stereo model \(i=Lg\),\allowbreak{} which claim is correct?
\begin{enumerate}[label=\textbf{\Alph*.}]
\item Only the albedo is ambiguous; the normal is always unique.
\item Twenty images guarantee a unique g.
\item The component of g perpendicular to that plane is unobservable.
\item Normalizing every light vector makes L full rank.
\end{enumerate}
\end{minipage}
\par\vspace{12pt}
\begin{minipage}{\linewidth}
\textbf{Q06 — Multiple choice}\par\smallskip
A photometric-stereo estimate has \(g=\begin{bmatrix}0&0&0\end{bmatrix}^{\mathsf T}\). What is the correct interpretation of \(n=\frac{g}{\lVert g\rVert}\)?
\begin{enumerate}[label=\textbf{\Alph*.}]
\item The normal must be \(\begin{bmatrix}0&0&1\end{bmatrix}\).
\item The surface must be infinitely far away.
\item The division produces a valid zero unit normal.
\item The normal is undefined; mask or otherwise handle the zero-reflectance/\allowbreak{}invalid pixel.
\end{enumerate}
\end{minipage}
\par\vspace{12pt}
\begin{minipage}{\linewidth}
\textbf{Q07 — Multiple choice}\par\smallskip
For normalized intensity \(0<r<1\),\allowbreak{} apply \(s=r^{\gamma}\) with \(0<\gamma<1\). What happens?
\begin{enumerate}[label=\textbf{\Alph*.}]
\item Every output intensity becomes 1.
\item The result is a linear contrast stretch.
\item Interior intensities increase; 0 and 1 remain fixed.
\item Interior intensities decrease; 0 and 1 remain fixed.
\end{enumerate}
\end{minipage}
\par\vspace{12pt}
\begin{minipage}{\linewidth}
\textbf{Q08 — Multiple choice}\par\smallskip
Why does a two-pass connected-component labelling algorithm store equivalences between provisional labels?
\begin{enumerate}[label=\textbf{\Alph*.}]
\item \(I_t\) estimates the intensity histogram before labelling.
\item Different provisional labels can later be found to belong to one connected region.
\item \(I_t\) needs to assign a different label to every foreground pixel.
\item \(I_t\) changes 4-connectivity into 8-connectivity automatically.
\end{enumerate}
\end{minipage}
\par\vspace{12pt}
\begin{minipage}{\linewidth}
\textbf{Q09 — Multiple choice}\par\smallskip
Foreground objects contain tiny isolated foreground specks and small background holes. Which sequence directly targets the two defects,\allowbreak{} assuming suitable structuring elements?
\begin{enumerate}[label=\textbf{\Alph*.}]
\item Closing,\allowbreak{} then another dilation.
\item Dilation,\allowbreak{} then another dilation.
\item Opening,\allowbreak{} then closing.
\item Erosion,\allowbreak{} then another erosion.
\end{enumerate}
\end{minipage}
\par\vspace{12pt}
\begin{minipage}{\linewidth}
\textbf{Q10 — Multiple choice}\par\smallskip
You implement a sliding dot product with an asymmetric 2D kernel K. What change converts this cross-correlation to mathematical convolution?
\begin{enumerate}[label=\textbf{\Alph*.}]
\item Flip K horizontally and vertically before the sliding dot product.
\item Divide K by its sum.
\item Negate K without changing its positions.
\item Transpose K only.
\end{enumerate}
\end{minipage}
\par\vspace{12pt}
\begin{minipage}{\linewidth}
\textbf{Q11 — Multiple choice}\par\smallskip
Two ideal Gaussian smoothing filters with standard deviations 3 and 4 are applied successively. Ignore sampling and truncation. What is the equivalent standard deviation?
\begin{enumerate}[label=\textbf{\Alph*.}]
\item 25
\item 12
\item 5
\item 7
\end{enumerate}
\end{minipage}
\par\vspace{12pt}
\begin{minipage}{\linewidth}
\textbf{Q12 — Multiple choice}\par\smallskip
A nonzero Fourier coefficient is \(a+\mathrm{i}b\). In NumPy's argument convention,\allowbreak{} which pair gives its magnitude and phase?
\begin{enumerate}[label=\textbf{\Alph*.}]
\item \(a+b,\;\operatorname{arctan2}(a,b)\)
\item \(\sqrt{a^2+b^2},\;\operatorname{arctan2}(b,a)\)
\item \(a^2+b^2,\;\arctan\!\left(\frac{b}{a}\right)\) for every a and b
\item \(\sqrt{a^2+b^2},\;\operatorname{arctan2}(a,b)\)
\end{enumerate}
\end{minipage}
\par\vspace{12pt}
\begin{minipage}{\linewidth}
\textbf{Q13 — Multiple choice}\par\smallskip
Under valid continuous-domain differentiation and convolution,\allowbreak{} which operation is equivalent to differentiating a Gaussian-smoothed image?
\begin{enumerate}[label=\textbf{\Alph*.}]
\item Convolve the image with the corresponding derivative of the Gaussian.
\item Threshold the image and then take a histogram.
\item Convolve the image with the square of the Gaussian.
\item Differentiate the Gaussian's standard deviation instead of a spatial coordinate.
\end{enumerate}
\end{minipage}
\par\vspace{12pt}
\begin{minipage}{\linewidth}
\textbf{Q14 — Multiple choice}\par\smallskip
After Canny non-maximum suppression,\allowbreak{} a chain of weak pixels connects to a strong pixel; another weak chain is disconnected from all strong pixels. Hysteresis should do what?
\begin{enumerate}[label=\textbf{\Alph*.}]
\item Keep every weak pixel because it exceeds the low threshold.
\item Discard every weak pixel because it is below the high threshold.
\item Keep whichever chain has the greatest number of pixels,\allowbreak{} regardless of connection.
\item Keep the connected weak chain and discard the disconnected chain.
\end{enumerate}
\end{minipage}
\par\vspace{12pt}
\begin{minipage}{\linewidth}
\textbf{Q15 — Multiple choice}\par\smallskip
A local structure tensor has two large eigenvalues. What does this indicate?
\begin{enumerate}[label=\textbf{\Alph*.}]
\item A guaranteed scale-invariant feature at every resolution.
\item Strong intensity variation in two independent directions: a good corner-like feature.
\item A flat patch.
\item A single straight edge with unconstrained tangential motion.
\end{enumerate}
\end{minipage}
\par\vspace{12pt}
\begin{minipage}{\linewidth}
\textbf{Q16 — Multiple choice}\par\smallskip
Why is \(y=mx+b\) with vertical-residual ordinary least squares unsuitable for an exactly vertical line?
\begin{enumerate}[label=\textbf{\Alph*.}]
\item \(I_t\) requires equal x and y measurement variances; unequal variances are the reason a vertical line is unrepresentable.
\item The model becomes valid once both coordinate columns are centred,\allowbreak{} because centring makes the slope finite.
\item Orthogonal-distance fitting handles vertical lines by forcing \(m=0\) in the same \(y=mx+b\) representation.
\item A vertical line cannot be represented by a finite m; an orthogonal-distance line model avoids this problem.
\end{enumerate}
\end{minipage}
\par\vspace{12pt}
\begin{minipage}{\linewidth}
\textbf{Q17 — Multiple choice}\par\smallskip
For fixed desired success probability and model sample size,\allowbreak{} what happens to the required RANSAC iteration count when the inlier fraction decreases?
\begin{enumerate}[label=\textbf{\Alph*.}]
\item \(I_t\) decreases because more points are rejected.
\item RANSAC becomes mathematically impossible as soon as inliers are below 50\%.
\item \(I_t\) increases,\allowbreak{} often sharply.
\item \(I_t\) is unchanged because RANSAC is robust.
\end{enumerate}
\end{minipage}
\par\vspace{12pt}
\begin{minipage}{\linewidth}
\textbf{Q18 — Multiple choice}\par\smallskip
In Lucas–Kanade,\allowbreak{} A has rows \(\begin{bmatrix}I_x&I_y\end{bmatrix}\). Which matrix also appears in corner analysis?
\begin{enumerate}[label=\textbf{\Alph*.}]
\item \(A^{\mathsf T}A\): the local second-moment/\allowbreak{}structure tensor.
\item \(AA^{\mathsf T}\),\allowbreak{} always a 2 by 2 matrix.
\item The intrinsic calibration matrix K.
\item The homography's last row.
\end{enumerate}
\end{minipage}
\par\vspace{12pt}
\begin{minipage}{\linewidth}
\textbf{Q19 — Multiple choice}\par\smallskip
A Lucas–Kanade window crosses the boundary between two independently moving objects. Which assumption is most directly violated?
\begin{enumerate}[label=\textbf{\Alph*.}]
\item The image contains any gradients at all.
\item The image has more than one colour channel.
\item The camera has a finite focal length.
\item A single common flow vector for all pixels in the window.
\end{enumerate}
\end{minipage}
\par\vspace{12pt}
\begin{minipage}{\linewidth}
\textbf{Q20 — Multiple choice}\par\smallskip
Which property is preserved by every nonsingular affine image transformation,\allowbreak{} but not by every projective transformation?
\begin{enumerate}[label=\textbf{\Alph*.}]
\item All Euclidean lengths.
\item All angles.
\item The location of the origin.
\item Parallelism of lines.
\end{enumerate}
\end{minipage}
\par\vspace{12pt}
\begin{minipage}{\linewidth}
\textbf{Q21 — Multiple choice}\par\smallskip
Two views of a static scene are taken by rotating a camera about exactly the same optical centre. The scene is not planar. Which model still relates the ideal undistorted images?
\begin{enumerate}[label=\textbf{\Alph*.}]
\item A unique depth-dependent map recovered by inverting the fundamental matrix.
\item No global model can exist unless the scene is planar.
\item One homography,\allowbreak{} independent of scene depth.
\item A 2D affine map for every possible rotation angle and calibration,\allowbreak{} without a projective denominator.
\end{enumerate}
\end{minipage}
\par\vspace{12pt}
\begin{minipage}{\linewidth}
\textbf{Q22 — Multiple choice}\par\smallskip
Why impose \(\lVert h\rVert=1\) when solving the homogeneous DLT problem \(\min_h\lVert Ah\rVert^2\) for a homography?
\begin{enumerate}[label=\textbf{\Alph*.}]
\item Because all nine entries of H must equal 1.
\item Without a scale constraint,\allowbreak{} \(h=0\) is a trivial minimizer and does not define a homography.
\item Because every image point must have unit Euclidean norm.
\item Because this fixes lens distortion.
\end{enumerate}
\end{minipage}
\par\vspace{12pt}
\begin{minipage}{\linewidth}
\textbf{Q23 — Multiple choice}\par\smallskip
A descriptor's nearest and second-nearest match distances are almost equal. Why does Lowe's ratio test tend to reject it?
\begin{enumerate}[label=\textbf{\Alph*.}]
\item Similar distances establish a reliable nearest match; the ratio test should therefore always keep it.
\item The match is ambiguous even if the nearest absolute distance is small.
\item The ratio test rejects it only because the nearest absolute distance must exceed a fixed universal cutoff.
\item \(I_t\) has already been proven to be a geometric outlier; descriptor distances directly establish reprojection error.
\end{enumerate}
\end{minipage}
\par\vspace{12pt}
\begin{minipage}{\linewidth}
\textbf{Q24 — Multiple choice}\par\smallskip
Two images have identical counts of visual words but different arrangements of those words. A plain Bag-of-Visual-Words histogram will do what?
\begin{enumerate}[label=\textbf{\Alph*.}]
\item Recover their exact layout from the counts.
\item Represent them identically,\allowbreak{} because it discards spatial arrangement.
\item Require a different number of histogram bins for each image.
\item Use the image labels as its histogram bin values.
\end{enumerate}
\end{minipage}
\par\vspace{12pt}
\begin{minipage}{\linewidth}
\textbf{Q25 — Multiple choice}\par\smallskip
You want a clean estimate of generalization for a BoVW classifier. Which protocol is best?
\begin{enumerate}[label=\textbf{\Alph*.}]
\item Choose k and vocabulary size using test accuracy,\allowbreak{} then report that accuracy.
\item Fit the vocabulary on all images,\allowbreak{} including the test set,\allowbreak{} to make evaluation easier.
\item Choose every setting using training accuracy alone.
\item Fit the vocabulary on training descriptors; select settings on validation data; evaluate once on untouched test data.
\end{enumerate}
\end{minipage}
\par\vspace{12pt}
\begin{minipage}{\linewidth}
\textbf{Q26 — Multiple choice}\par\smallskip
A classifier predicts the majority class for every example and achieves 99\% accuracy on a 99:1 class split. What can its minority-class recall be?
\begin{enumerate}[label=\textbf{\Alph*.}]
\item 0
\item 1
\item 0.99
\item \(I_t\) must equal minority-class precision and be positive.
\end{enumerate}
\end{minipage}
\par\vspace{12pt}
\begin{minipage}{\linewidth}
\textbf{Q27 — Multiple choice}\par\smallskip
Ignoring boundaries,\allowbreak{} a stride-1 convolution with shared weights most directly provides which property?
\begin{enumerate}[label=\textbf{\Alph*.}]
\item Exact translation invariance: the whole feature map never changes after a shift.
\item A separate learned kernel for each image location.
\item Exact invariance to every rotation and scale change.
\item Translation equivariance: shifting the input shifts the feature map.
\end{enumerate}
\end{minipage}
\par\vspace{12pt}
\begin{minipage}{\linewidth}
\textbf{Q28 — Multiple choice}\par\smallskip
What is the main proposal-generation difference between Fast R-CNN and Faster R-CNN?
\begin{enumerate}[label=\textbf{\Alph*.}]
\item Fast R-CNN predicts pixel masks instead of boxes.
\item Fast R-CNN has an RPN and Faster R-CNN removes it.
\item Faster R-CNN uses a learned Region Proposal Network; Fast R-CNN normally relies on external proposals such as Selective Search.
\item Faster R-CNN runs a separate backbone independently on every raw crop.
\end{enumerate}
\end{minipage}
\par\vspace{12pt}
\begin{minipage}{\linewidth}
\textbf{Q29 — Multiple choice}\par\smallskip
Which statement correctly distinguishes the usual skip operations of ResNet and U-Net?
\begin{enumerate}[label=\textbf{\Alph*.}]
\item ResNet adds compatible activations; U-Net concatenates encoder and decoder features across matching scales.
\item Both use elementwise addition with mismatched channel counts; the addition automatically supplies the missing channels.
\item Both normally concatenate along the channel axis,\allowbreak{} so every skip doubles the number of channels.
\item ResNet concatenates encoder/\allowbreak{}decoder features; U-Net adds an identity tensor to each residual branch.
\end{enumerate}
\end{minipage}
\par\vspace{12pt}
\begin{minipage}{\linewidth}
\textbf{Q30 — Multiple choice}\par\smallskip
Given a homogeneous image point x in image 1 and fundamental matrix F,\allowbreak{} with \(Fx\) nonzero,\allowbreak{} what does \(Fx\) specify in image 2?
\begin{enumerate}[label=\textbf{\Alph*.}]
\item The unique matching point,\allowbreak{} including its depth.
\item An epipolar line on which the corresponding point must lie.
\item A viewing-ray depth that becomes metric once F is normalized to unit Frobenius norm.
\item The matching image point after dehomogenizing \(Fx\) as if it were a homography output.
\end{enumerate}
\end{minipage}
\par\vspace{12pt}
\clearpage\section*{Part B — Numerical values}
\begin{minipage}{\linewidth}
\textbf{Q31 — Numerical}\par\smallskip
Use column vectors and the virtual image plane in front of the camera. The camera centre is \(C=\begin{bmatrix}1&2&0\end{bmatrix}^{\mathsf T}\),\allowbreak{} \(X_w=\begin{bmatrix}3&5&4\end{bmatrix}^{\mathsf T}\),\allowbreak{} and \(X_c=R(X_w-C)\).


\begin{Code}
import numpy as np
C = np.array([1,2,0])
Xw = np.array([3,5,4])
R = np.array([[0,-1,0], [1,0,0], [0,0,1]])
K = np.array([[800,20,320], [0,600,240], [0,0,1]])
\end{Code}


Projection is \(q=KX_c\),\allowbreak{} followed by \(u=\frac{q_0}{q_2}\). Report u in pixels.
\par\medskip\textbf{Answer:} \rule{55mm}{0.4pt}
\end{minipage}
\par\vspace{12pt}
\begin{minipage}{\linewidth}
\textbf{Q32 — Numerical}\par\smallskip
For one pixel,\allowbreak{} use \(i=Lg\),\allowbreak{} where light strength and camera gain are already absorbed into L. Compute the least-squares estimate,\allowbreak{} then report albedo \(\rho=\lVert g\rVert\).


\begin{Code}
import numpy as np
L = np.array([[0,0,1], [.8,0,.6], [-.8,0,.6],
              [0,.8,.6], [0,-.8,.6]])
i = np.array([.70,.64,.16,.44,.36])
\end{Code}


Provided routine: \texttt{g =\allowbreak{} np.linalg.lstsq(L,\allowbreak{} i,\allowbreak{} rcond=\allowbreak{}None)[0]} minimizes the sum of squared residuals. \texttt{np.linalg.norm(g)} returns the Euclidean norm.
\par\medskip\textbf{Answer:} \rule{55mm}{0.4pt}
\end{minipage}
\par\vspace{12pt}
\begin{minipage}{\linewidth}
\textbf{Q33 — Numerical}\par\smallskip
Use the height-field convention \(n\propto\begin{bmatrix}-p&-q&1\end{bmatrix}^{\mathsf T}\),\allowbreak{} where \(p=\frac{\partial z}{\partial x}\) and \(q=\frac{\partial z}{\partial y}\). A plane has constant normal proportional to \(\begin{bmatrix}2&-1&2\end{bmatrix}\) and \(z(0,0)=5\). Report \(z(3,4)\),\allowbreak{} using unit spatial spacing.
\par\medskip\textbf{Answer:} \rule{55mm}{0.4pt}
\end{minipage}
\par\vspace{12pt}
\begin{minipage}{\linewidth}
\textbf{Q34 — Numerical}\par\smallskip
A 2-bit image has histogram counts \(\begin{bmatrix}1&3&2&2\end{bmatrix}\) for levels \(\begin{bmatrix}0&1&2&3\end{bmatrix}\). Use this specific equalization rule: \(T(r)=\lfloor3\operatorname{CDF}(r)\rfloor\),\allowbreak{} with no subtraction of \(\operatorname{CDF}_{\min}\). Report \(T(1)\) as an exact integer.
\par\medskip\textbf{Answer:} \rule{55mm}{0.4pt}
\end{minipage}
\par\vspace{12pt}
\begin{minipage}{\linewidth}
\textbf{Q35 — Numerical}\par\smallskip
Close the following 1D binary row with the centred structuring element \(\begin{bmatrix}1&1&1\end{bmatrix}\). Closing means dilation followed by erosion. At each stage retain exactly 11 entries and treat all out-of-range values as 0. Report the number of foreground entries after closing.


\[
\begin{bmatrix}0&1&1&0&1&1&0&0&1&0&0\end{bmatrix}
\]

\par\medskip\textbf{Answer:} \rule{55mm}{0.4pt}
\end{minipage}
\par\vspace{12pt}
\begin{minipage}{\linewidth}
\textbf{Q36 — Numerical}\par\smallskip
For the 3 by 3 patch I and kernel K below,\allowbreak{} report the single valid output of mathematical convolution,\allowbreak{} not cross-correlation. The kernel is flipped on both axes. Report an exact integer.


\[
\begin{aligned}I&=\begin{bmatrix}1&2&0\\3&-1&4\\2&1&5\end{bmatrix}\\[6pt]K&=\begin{bmatrix}1&0&2\\0&-1&0\\3&0&-2\end{bmatrix}\end{aligned}
\]

\par\medskip\textbf{Answer:} \rule{55mm}{0.4pt}
\end{minipage}
\par\vspace{12pt}
\begin{minipage}{\linewidth}
\textbf{Q37 — Numerical}\par\smallskip
Use the unnormalized DFT \(F_k=\sum_{n=0}^{3}f_n\exp\!\left(-\frac{2\pi\mathrm{i}kn}{4}\right)\). For \(f=\begin{bmatrix}1&2&3&4\end{bmatrix}\),\allowbreak{} report the magnitude of \(F_1\). You can use \(\exp(-\mathrm{i}\pi/2)=-\mathrm{i}\) and \(\exp(-\mathrm{i}\pi)=-1\).
\par\medskip\textbf{Answer:} \rule{55mm}{0.4pt}
\end{minipage}
\par\vspace{12pt}
\begin{minipage}{\linewidth}
\textbf{Q38 — Numerical}\par\smallskip
A patch has structure tensor
\[M=\begin{bmatrix}9&3\\3&9\end{bmatrix}.\]
 The Harris response is \(R=\det(M)-k\operatorname{tr}(M)^2\),\allowbreak{} with \(k=0.04\). Report R.
\par\medskip\textbf{Answer:} \rule{55mm}{0.4pt}
\end{minipage}
\par\vspace{12pt}
\begin{minipage}{\linewidth}
\textbf{Q39 — Numerical}\par\smallskip
Use \(I_x u+I_y v=-I_t\) and one constant flow in the window. Compute the least-squares flow and report u.


\begin{Code}
import numpy as np
A = np.array([[1,0], [0,1], [1,1], [2,0]], dtype=float)
It = np.array([-2,1,-1,-5], dtype=float)
\end{Code}


Provided routine: \texttt{np.linalg.lstsq(A,\allowbreak{} -It,\allowbreak{} rcond=\allowbreak{}None)[0]} minimizes squared residuals.
\par\medskip\textbf{Answer:} \rule{55mm}{0.4pt}
\end{minipage}
\par\vspace{12pt}
\begin{minipage}{\linewidth}
\textbf{Q40 — Numerical}\par\smallskip
A homography RANSAC uses samples of 4 correspondences,\allowbreak{} each independently an inlier with probability \(w=0.6\). Report the minimum integer iteration count for success probability at least \(p=0.99\). Use \(N=\left\lceil\frac{\log(1-p)}{\log(1-w^4)}\right\rceil\). Assume all-inlier samples are nondegenerate.
\par\medskip\textbf{Answer:} \rule{55mm}{0.4pt}
\end{minipage}
\par\vspace{12pt}
\begin{minipage}{\linewidth}
\textbf{Q41 — Numerical}\par\smallskip
A homography maps source to target with
\[H=\begin{bmatrix}2&0&1\\0&1&-2\\0.1&0&1\end{bmatrix}.\]
 Use column-vector homogeneous coordinates. Report the target x coordinate of source point (5,\allowbreak{}4).
\par\medskip\textbf{Answer:} \rule{55mm}{0.4pt}
\end{minipage}
\par\vspace{12pt}
\begin{minipage}{\linewidth}
\textbf{Q42 — Numerical}\par\smallskip
Backward warping maps a target pixel to source coordinates \((x,y)=(0.25,0.75)\). Source intensities are \(I_{0,0}=10\),\allowbreak{} \(I_{0,1}=30\),\allowbreak{} \(I_{1,0}=50\),\allowbreak{} \(I_{1,1}=90\). Arrays use \(I_{y,x}\). Report the bilinear interpolated intensity. The weights are \((1-x)(1-y),\;x(1-y),\;(1-x)y\text{ and }xy\),\allowbreak{} respectively.
\par\medskip\textbf{Answer:} \rule{55mm}{0.4pt}
\end{minipage}
\par\vspace{12pt}
\begin{minipage}{\linewidth}
\textbf{Q43 — Numerical}\par\smallskip
Rows of the following confusion matrix are actual classes A,\allowbreak{}B,\allowbreak{}C; columns are predicted classes A,\allowbreak{}B,\allowbreak{}C. Report the F1 score for class B as a fraction,\allowbreak{} not a percentage.


\[
\begin{bmatrix}18&2&0\\3&12&5\\1&4&15\end{bmatrix}
\]

\par\medskip\textbf{Answer:} \rule{55mm}{0.4pt}
\end{minipage}
\par\vspace{12pt}
\begin{minipage}{\linewidth}
\textbf{Q44 — Numerical}\par\smallskip
A retrieval query has 4 relevant database images in total. In its returned ranking,\allowbreak{} relevant images occur at ranks 1,\allowbreak{}3,\allowbreak{}5; the fourth is not returned. Use uninterpolated
\[\mathrm{AP}=\frac{\text{sum of precision at each relevant rank}}{\text{total relevant images}}.\]
 Report AP as a fraction.
\par\medskip\textbf{Answer:} \rule{55mm}{0.4pt}
\end{minipage}
\par\vspace{12pt}
\begin{minipage}{\linewidth}
\textbf{Q45 — Numerical}\par\smallskip
Boxes use continuous coordinates \([x_{\min},y_{\min},x_{\max},y_{\max}]\),\allowbreak{} so \(\text{width}=x_{\max}-x_{\min}\),\allowbreak{} without a +1 correction. Report IoU for \(A=\begin{bmatrix}0&0&6&4\end{bmatrix}\) and \(B=\begin{bmatrix}2&1&8&5\end{bmatrix}\).
\par\medskip\textbf{Answer:} \rule{55mm}{0.4pt}
\end{minipage}
\par\vspace{12pt}
\begin{minipage}{\linewidth}
\textbf{Q46 — Numerical}\par\smallskip
A standard 2D convolution takes an \(H\times W\times C_{\mathrm{in}}\) input of \(31\times29\times3\). Kernel height/\allowbreak{}width =\allowbreak{} \(5\times3\),\allowbreak{} stride =\allowbreak{} \(2\times2\),\allowbreak{} padding =\allowbreak{} \(1\times0\) on EACH side,\allowbreak{} and \(C_{\mathrm{out}}=8\). No dilation. Report the total number of scalar outputs,\allowbreak{} excluding the batch dimension. Use
\[\text{out}=\left\lfloor\frac{\text{in}+2\cdot\text{pad}-\text{kernel}}{\text{stride}}\right\rfloor+1\text{ per axis}.\]

\par\medskip\textbf{Answer:} \rule{55mm}{0.4pt}
\end{minipage}
\par\vspace{12pt}
\begin{minipage}{\linewidth}
\textbf{Q47 — Numerical}\par\smallskip
For the convolution in Q46,\allowbreak{} report the number of trainable parameters INCLUDING one bias per output channel. \(I_t\) is a dense convolution across all 3 input channels; no groups.
\par\medskip\textbf{Answer:} \rule{55mm}{0.4pt}
\end{minipage}
\par\vspace{12pt}
\begin{minipage}{\linewidth}
\textbf{Q48 — Numerical}\par\smallskip
Flatten a \(7\times7\times16\) feature map and feed it to a fully connected layer with 10 output units. Include one bias per output unit. Report its trainable parameter count.
\par\medskip\textbf{Answer:} \rule{55mm}{0.4pt}
\end{minipage}
\par\vspace{12pt}
\begin{minipage}{\linewidth}
\textbf{Q49 — Numerical}\par\smallskip
A network applies: a \(3\times3\) convolution with stride 1; a \(2\times2\) pooling layer with stride 2; a \(3\times3\) convolution with stride 1. Report the side length of one final unit's theoretical receptive field in original input pixels. Ignore boundaries. Provided recurrence: \(r_{\mathrm{new}}=r_{\mathrm{old}}+(\text{kernel}-1)j_{\mathrm{old}}\); \(j_{\mathrm{new}}=j_{\mathrm{old}}\cdot\text{stride}\); start \(r=j=1\).
\par\medskip\textbf{Answer:} \rule{55mm}{0.4pt}
\end{minipage}
\par\vspace{12pt}
\begin{minipage}{\linewidth}
\textbf{Q50 — Numerical}\par\smallskip
For \(x_2^{\mathsf T}Fx_1=0\),\allowbreak{} let
\[F=\begin{bmatrix}0&0&0\\0&0&-1\\0&1&0\end{bmatrix},\]
 \(x_{1}=\begin{bmatrix}12&7&1\end{bmatrix}^{\mathsf T}\) and candidate \(x_{2}=\begin{bmatrix}9&9&1\end{bmatrix}^{\mathsf T}\). Report the perpendicular distance in pixels from x2 to its epipolar line \(l_2=Fx_1\). Provided formula: \(\text{distance}=\frac{\lvert ax+by+c\rvert}{\sqrt{a^2+b^2}}\),\allowbreak{} for line \([a,b,c]\).
\par\medskip\textbf{Answer:} \rule{55mm}{0.4pt}
\end{minipage}
\par\vspace{12pt}

\end{document}
